\documentclass[11pt]{article}
\usepackage[T1]{fontenc}
\usepackage[margin=27mm]{geometry}
\usepackage{amsmath,amssymb,amsthm,booktabs,longtable,hyperref}
\newcommand{\Q}{\mathbb Q}
\newcommand{\Fp}{\mathbb F_p}
\newcommand{\PP}{\mathbb P}
\newcommand{\Pic}{\operatorname{Pic}}
\newcommand{\gon}{\operatorname{gon}}
\newtheorem{theorem}{Theorem}
\newtheorem{lemma}[theorem]{Lemma}
\newtheorem{proposition}[theorem]{Proposition}
\theoremstyle{remark}\newtheorem{remark}[theorem]{Remark}
\title{Intermediate modular curves with infinitely many quintic points}
\author{Proof and computational certificates}
\date{6 September 2026}
\begin{document}
\maketitle
\begin{abstract}
We classify the intermediate modular curves with infinitely many closed
points of degree exactly five over $\Q$. There are 43 positive strict
subgroup-defined pairs, each admitting a rational function of degree five.
We separate the short mathematical arguments from the exact computations
they require. Complete case tables, source checks, and freshly replayed
computational certificates accompany the proof.
\end{abstract}

\section{Statement and conventions}
For $\{\pm1\}\subseteq\Delta\subseteq(\mathbb Z/N\mathbb Z)^\times$, set
\[
 \Gamma_\Delta(N)=\left\{\begin{pmatrix}a&b\\c&d\end{pmatrix}
 \in\mathrm{SL}_2(\mathbb Z):c\equiv0\pmod N,\ a\bmod N\in\Delta\right\}.
\]
The smooth projective modular curve attached to this subgroup is denoted
$X_\Delta(N)$. We keep the pairs $(N,\Delta)$, rather than identifying
curves that happen to be isomorphic. A pair is \emph{strict} if both subgroup
inclusions above are strict. A quintic point means a closed point with
residue field of degree exactly five, not a point of degree at most five.

\begin{theorem}[Classification]
The strict intermediate curves with infinitely many quintic points are
the 43 pairs in Table~\ref{tab:positive}. Each admits a $\Q$-rational
function of degree exactly five. Every other strict intermediate curve
has only finitely many quintic points.

Including endpoints, $X_1(N)$ is positive precisely for
$N\in\{1,\ldots,25\}\setminus\{23\}$. The positive $X_0$ levels are
\[
 \{1,\ldots,45\}\cup\{47,\ldots,58\}\cup
 \{61,63,64,65,67,68,72,73,75,80,81,91,109,121,125\}.
\]
\end{theorem}

The endpoint assertions are inputs to the strict classification. The
$X_1$ assertion comes from Derickx--Sutherland \cite{DS}; the $X_0$ assertion uses
the published classification \cite[Theorem 1.10]{DHJO} together with the
thirty-level proof in \texttt{../A\_x0\_quintic/quintic\_points.pdf}
(Theorem A of the paper). Its three routes use rank-zero lifting, quotient
class sieves and local rigidity of reducible divisor families. Its 36
computational certificates are in \texttt{../A\_x0\_quintic/}.

For completeness, the $X_1$ endpoint statement concerns exact degree five.
The cyclic torsion orders in \cite[Theorem 1.1]{DS} are precisely
$1\le N\le25$, $N\ne23$. At genus zero, the rational cusp gives
$X_1(N)\simeq\PP^1$. At positive genus in this range, rank zero from
\cite[Theorem 4.1]{DS} makes rational points finite, so the infinitely
many marked elliptic curves over quintic fields give genuine quintic
points. Conversely, infinitely many quintic points give infinitely many
marked elliptic curves; uniform boundedness and finite fibers of the
$j$-map force a torsion group occurring infinitely often. Its cyclic
subgroup order must lie in the same list. Finally,
\cite[Corollary 2.4]{DS} gives exact degree-five functions at the positive
rank-zero endpoints as well.

\begin{table}[ht]
\centering
\begin{tabular}{lr}\toprule
Family of arguments & Strict pairs through level 191\\\midrule
Projection and spectral gonality &829\\
Rank-zero Jacobians &82\\
Positive-rank density and elliptic families &57\\
Degree-five function constructions &43\\\bottomrule
\end{tabular}
\caption{The four main groups of proof assignments.}
\end{table}

\section{Three elementary reductions}
\begin{lemma}[A function produces quintic points]
If a geometrically integral curve $C/\Q$ admits a degree-five map to
$\PP^1$, then it has infinitely many quintic points.
\end{lemma}
\begin{proof}
Hilbert irreducibility supplies infinitely many rational fibers whose
degree-five coordinate algebra is a field. After omitting branch values,
these fibers are closed points of degree five.
\end{proof}

\begin{lemma}[Projection]
Let $C\to D$ be finite over $\Q$. If $D$ has finitely many rational
points and finitely many quintic points, then $C$ has finitely many
quintic points.
\end{lemma}
\begin{proof}
The residue degree of the image divides five, hence is one or five.
There are finitely many such images and each fiber has finite support.
\end{proof}
The rational-point hypothesis cannot be dropped. In applications it
follows from $g(D)>1$ and Faltings' theorem \cite{Faltings}.
A negative $X_0$ target cannot
have genus at most two, since its rational cusp and the low-genus
construction below would give infinitely many quintic points.
Forgetting maps are defined
over $\Q$: for $M\mid N$, reduction of $\Delta_N$ into $\Delta_M$
permits $(E,P)\mapsto(E,(N/M)P)$. Same-level subgroup inclusions are
special cases. Iteration starts only from proved negatives, so it creates
no circular dependencies.

\begin{lemma}[Rank zero]
Suppose $C(\Q)\ne\varnothing$ and $J(C)(\Q)$ has rank zero.
If $C$ has no rational function of degree five, it has finitely many
quintic points.
\end{lemma}
\begin{proof}
A rational base point identifies $\Pic^5(C)(\Q)$ with the finite group
$J(C)(\Q)$. Let $D$ be an irreducible effective quintic divisor. If
$h^0(D)>1$, a nonconstant rational function in $L(D)$ has a nonzero pole
divisor bounded by $D$. Since $D$ is one closed point with coefficient
one, that pole divisor must equal $D$. This is a degree-five function,
a contradiction. Otherwise $D$ is the unique effective representative
of its class. There are only finitely many classes.
\end{proof}
\begin{remark}
The rank-zero assumption is essential to this argument. It will not be
used implicitly in the positive-rank cases.
\end{remark}

\section{Enumeration and rational models}
Primitive rows $(c,d)$ modulo $N$, modulo multiplication by $\Delta$,
parametrize $\Gamma_\Delta(N)\backslash\mathrm{SL}_2(\mathbb Z)$.
The scalar action is free: if $u(c,d)=(c,d)$, primitivity implies $u=1$.
Thus the index is
\[
 \mu=\frac{N^2}{|\Delta|}\prod_{\ell\mid N}(1-\ell^{-2}).
\]
Right multiplication by the standard generators gives permutations
$(c,d)\mapsto(d,-c)$ and $(c,d)\mapsto(c,c+d)$. Fixed points of the
order-two and order-three permutations count the elliptic points;
cycles of the latter translation count cusps. Therefore
\[
 g=1+\frac\mu{12}-\frac{e_2}4-\frac{e_3}3-\frac{c_\infty}2.
\]
Enumerating subgroups by successively adjoining every unit is complete:
every subgroup has a finite generating set. Full element sets, not their
orders or genera, are used as identifiers. A second computation with
Sage's abstract unit groups and \texttt{GammaH} checks all 1388 pairs
through level 191, of which 1011 are strict.

The rational structure deserves explicit attention. Rational Fourier
expansions at the Tate cusp describe the model parametrizing
$\mu_N$-markings, not automatically the analytically identical
point-marked model \cite[proof of Lemma 6.1.6]{CES}. Quotienting an elliptic curve by its marked subgroup
and dualizing the isogeny interchanges a $\mu_N$-marking and a
$\mathbb Z/N\mathbb Z$-marking. These inverse constructions are over
$\Q$ and interchange diamond $a$ with diamond $a^{-1}$. Since $\Delta$
is closed under inversion, they descend to a $\Q$-isomorphism of the
two intermediate models. In particular all these curves have a rational
Tate cusp after this identification.

For genus at least three, the full weight-two cusp space supplies
canonical differentials. A degree-$d$ relation has weight $2d$ and is
certified by vanishing through $\lfloor2d\mu/12\rfloor$. Extra Fourier
coefficients are checked but do not replace this bound. The quadratic
product rank distinguishes the hyperelliptic canonical image from the
nonhyperelliptic case. In the latter case the known canonical image is
embedded, of degree $2g-2$ and genus $g$. Exact containment in the proposed
scheme, followed by equality of Hilbert polynomials, identifies the whole
scheme: the quotient ideal sheaf has zero Hilbert polynomial and hence
vanishes. Merely matching genera would not suffice.

Before a reduction argument, the displayed equations must pass a separate
test at the chosen prime. Their special fiber is required to have the same
Hilbert polynomial as the generic canonical curve, to be reduced and
irreducible, to have normalization genus equal to its arithmetic genus,
and to have a rational point. Here is why these tests suffice. The flat
closure of the generic curve is contained in the displayed integral scheme.
Its special fiber has the generic Hilbert polynomial. Equality with the
displayed special-fiber Hilbert polynomial therefore identifies these two
special fibers by the same quotient-sheaf argument as above. Equality of
normalization and arithmetic genera makes the integral special curve
normal, hence smooth over the perfect residue field. Its rational point
forces the constant field to be the residue field, giving geometric
integrality. The flat closure is consequently a smooth proper model near
the special fiber. This is the model on which specialization is used;
the condition $p\nmid N$ alone is not a check on chosen coordinates.

\section{Positive cases}
If $g\le2$ and $P\in C(\Q)$, Riemann--Roch gives
$h^0(5P)=6-g$ and $h^0(4P)=5-g$. Thus a rational function in
$L(5P)\setminus L(4P)$ has pole divisor exactly $5P$, giving a
degree-five map.

In genus three, the echelon pivot exponents of the Fourier basis are
the Weierstrass gaps at the Tate cusp: a form of order $a$ yields a
differential of order $a-1$. If five is not a gap, then
$h^0(5P)>h^0(4P)$, giving a function with pole divisor $5P$.

In genus four, nonhyperellipticity makes the canonical bundle separate
all length-two subschemes. Therefore $K-P$ is base-point-free, including
at $P$, and has degree five. A general pair of sections again gives
the desired map. Nonhyperellipticity is checked by quadratic product
rank nine rather than seven.

The remaining twelve new positive cases use explicitly verified degree-five
pole divisors on certified models. The four published cases are exactly
those in Orli\'c's Theorems 1.4--1.5 \cite{Orlic}.
Fresh model generation and the degree-five witness checks have been replayed
for all twelve new cases. The fourteen canonical low-genus cases have also
been replayed; the thirteen genus-at-most-two cases use the uniform argument
above. The certificates are listed in \texttt{README.md}.

\clearpage
\begin{center}
\textbf{Positive strict pairs.}\quad
L: genus at most two; C: canonical low-genus argument;
W: explicit witness; P: published pentagonality.
\end{center}
\begingroup\small
\begin{longtable}{rrrl}
\caption{Positive strict intermediate pairs.}\label{tab:positive}\\
\toprule
$N$ & Generators of $\Delta$ & $g$ & Proof \\
\midrule\endfirsthead
\toprule $N$ & Generators of $\Delta$ & $g$ & Proof \\
\midrule\endhead
13 & $\langle -1, 3\rangle$ & 0 & L \\
13 & $\langle -1, 5\rangle$ & 0 & L \\
15 & $\langle -1, 4\rangle$ & 1 & L \\
16 & $\langle -1, 7\rangle$ & 0 & L \\
17 & $\langle -1, 2\rangle$ & 1 & L \\
17 & $\langle -1, 4\rangle$ & 1 & L \\
19 & $\langle -1, 7\rangle$ & 1 & L \\
20 & $\langle -1, 9\rangle$ & 1 & L \\
21 & $\langle -1, 4\rangle$ & 1 & L \\
21 & $\langle -1, 8\rangle$ & 3 & C \\
24 & $\langle -1, 5\rangle$ & 3 & C \\
24 & $\langle -1, 7\rangle$ & 3 & C \\
24 & $\langle -1, 11\rangle$ & 1 & L \\
25 & $\langle -1, 4\rangle$ & 0 & L \\
25 & $\langle -1, 7\rangle$ & 4 & C \\
26 & $\langle -1, 3\rangle$ & 4 & C \\
26 & $\langle -1, 5\rangle$ & 4 & C \\
27 & $\langle -1, 8\rangle$ & 1 & L \\
28 & $\langle -1, 3\rangle$ & 4 & C \\
28 & $\langle -1, 13\rangle$ & 4 & C \\
29 & $\langle -1, 4\rangle$ & 4 & C \\
30 & $\langle -1, 11\rangle$ & 5 & W \\
31 & $\langle -1, 2\rangle$ & 6 & P \\
31 & $\langle -1, 5\rangle$ & 6 & P \\
32 & $\langle -1, 7\rangle$ & 1 & L \\
32 & $\langle -1, 15\rangle$ & 5 & W \\
33 & $\langle -1, 2\rangle$ & 5 & W \\
34 & $\langle -1, 9\rangle$ & 5 & W \\
35 & $\langle -1, 4\rangle$ & 5 & W \\
36 & $\langle -1, 11\rangle$ & 3 & C \\
37 & $\langle -1, 3\rangle$ & 4 & C \\
37 & $\langle -1, 8\rangle$ & 4 & C \\
39 & $\langle -1, 4\rangle$ & 5 & W \\
40 & $\langle -1, 7\rangle$ & 7 & W \\
40 & $\langle -1, 9, 11\rangle$ & 5 & W \\
41 & $\langle -1, 2\rangle$ & 5 & W \\
44 & $\langle -1, 5\rangle$ & 8 & P \\
45 & $\langle -1, 4\rangle$ & 5 & W \\
48 & $\langle -1, 11\rangle$ & 5 & W \\
49 & $\langle -1, 6\rangle$ & 3 & C \\
50 & $\langle -1, 9\rangle$ & 4 & C \\
64 & $\langle -1, 7\rangle$ & 5 & W \\
125 & $\langle -1, 4\rangle$ & 16 & P \\
\bottomrule
\end{longtable}

\endgroup

\section{Uniform negative arguments}
Projection to the completed $X_0$ classification excludes every level
$N\ge192$. It excludes 807 strict pairs below this boundary. Seven more
pairs are excluded by other forgetting maps. The spectral bound
$\gon_{\mathbb C}C\ge119\mu/12000$ \cite[Theorem 3.16]{Orlic}, together with Frey's necessary bound
\cite[Proposition 2.2]{DS}
for quintic density, excludes another fifteen pairs: the implementation
uses the strict integer inequality $119\mu>120000$.

For rank-zero curves, it remains to exclude degree-five functions.
The published classifications of hyperelliptic, trigonal, tetragonal
and pentagonal intermediate curves do this whenever $\gon_{\Q}C\ge6$.
The subgroup lists are expanded as element sets before comparison;
genus or group order alone is not a sufficient identifier.
Nonvanishing central values establish rank zero by the modular
nonvanishing theorem. For the nontrivial-nebentypus factors occurring here,
we use Kato's Corollary 14.3, as applied to modular Jacobians in
Derickx--Sutherland, Section 4 \cite{Kato,DS}; citing only the
trivial-nebentypus version would not justify all intermediate cases.
The Magma overload used here factors the modular
abelian variety and tests exact modular-symbol $L$-ratios. It is not a
numerical approximation to $L(1)$. Rank zero of $J_1(N)$ also descends
to its intermediate quotients up to isogeny.

Two rank-zero exceptions, of genera seven and nine at levels 35 and 55,
are geometrically bielliptic \cite[Theorem 0.2 and Table 6]{JKS}.
A degree-two elliptic map and a degree-five
rational map generate the function field, since their degrees are coprime.
Castelnuovo--Severi would give $g\le2+(2-1)(5-1)=6$. Thus no quintic
function exists even geometrically. A tetragonal curve is not excluded
just because its gonality happens to equal four.

\section{Exhaustive rational divisor-class sieves}
Let $C$ have good reduction at an odd prime $p$ and a rational point $P$.
Suppose rank zero has been proved and an explicit subgroup of rational
divisor classes is proved to exhaust $J(C)(\Q)$. Reduction on rational
torsion over the unramified field $\Q_p$ is injective. Enumerate every
effective degree-five divisor on $C_{\Fp}$, using all closed places of
degrees one through five and all partitions of five. Keep the classes
with $h^0\ge2$ that belong to the known rational subgroup, shifted by
$5P$. This includes the specialization of every rational quintic pencil,
by upper semicontinuity. Over a finite field, an effective rational
line-bundle class has a rational effective representative, so enumeration
of rational effective divisors omits no relevant class.

Lift each surviving class to a rational divisor and compute its full
Riemann--Roch space. A bounded search may find a convenient representative,
but every surviving class must be checked. Injectivity of reduction
identifies the lift with the intended rational class. If its section
space has dimension at most one, it is not a pencil. If the dimension
is two and the complete system has a base point, it cannot give a
degree-five map. These checks exclude every rational quintic function;
the rank-zero lemma then excludes quintic density.

The assertion that the rational subgroup is complete is crucial.
At level 40, a full integral Manin-symbol lattice, intersected with the
cusp-boundary kernel, gives absolute integral homology. On its rational
quotient by this lattice, rational torsion is fixed by complex conjugation
and annihilated, away from the residue characteristic, by
\[
 1-T_p+p\langle p\rangle=1-T_p+T_p^2-T_{p^2}.
\]
Smith normal forms give rigorous upper bounds, with each residue prime
omitted when bounding its own primary part. Their combination with
Frobenius-order bounds gives 4032. Explicit rational point differences
generate 4032 classes on reduction, proving equality. Twelve special
classes remain, and all twelve are checked over $\Q$.

For the two level-48 cases, use the published description of
$J_1(48)(\Q)_{\rm tors}$ by rational cusp-orbit divisors
\cite[Theorem 4.13]{DEH}. Their embedding
is recomputed in full integral homology. Transfer for the degree-four
forgetting map is obtained by summing lifted paths over all sheets;
it induces pullback on Jacobians. Every rational torsion class on the
target must transfer into the computed rational cuspidal group upstairs.
This restricts the target Hecke bounds to 2304 and 12800, respectively.
Explicit closed-cusp classes attain both bounds. The two finite-field
sieves then leave four and 24 classes, all checked rationally. Pullback
need not be injective: taking its inverse image is the correct restriction.

Characterwise rational decompositions are used only for isogeny invariants,
not as integral homology lattices. For a rational-point divisor $D$, the
final Riemann--Roch computations can instead use canonical hyperplanes:
the saturated ideal $I_C+I(P)^m$ cuts out $mP$ on the smooth curve.
Its linear part is $H^0(K-mP)$. Intersect these spaces and apply
Riemann--Roch. Equalities $h^0(D-P)=h^0(D)$ certify base points.

\section{A formal lifting lemma}
The same local argument underlies the two remaining computational methods.
It is stronger than a dimension count on the special fiber.
\begin{lemma}\label{lem:formal}
Let $R$ be a complete discrete valuation ring. Let $\mathcal P/R$ be
smooth of relative dimension $g$, and $\mathcal T\subseteq\mathcal P$
a closed subscheme. Suppose a smooth formal $r$-dimensional branch
through a residue point maps into $\mathcal T$ and is a formal immersion
into $\mathcal P$. If the relative tangent space of $\mathcal T$ at
that residue point has dimension at most $r$, then the completion of
$\mathcal T$ there is exactly that branch.
\end{lemma}
\begin{proof}
Write the completed ambient ring as $R[[x_1,\ldots,x_g]]$. The tangent
bound supplies $g-r$ equations in the ideal of $\mathcal T$ with
independent transverse linear terms modulo the maximal ideal of $R$.
The formal implicit function theorem eliminates these variables and
leaves $R[[t_1,\ldots,t_r]]$. The contained branch has the same tangent
dimension, so its map into this ring is an isomorphism by the formal
inverse function theorem. Every remaining equation vanishes on that
branch and hence is zero. In particular, no additional vertical or
infinitesimal components survive.
\end{proof}

\subsection{Petri lifting for rank-zero curves}
Suppose explicit base-point-free quartic bundles $A_i$ have $h^0=2$
over $\Q$ and after good reduction. Suppose an exhaustive computation
shows that every class in $W^1_5(\Fp)$ is some $A_i+P$, with
$P\in C(\Fp)$, has exactly two sections, and has Petri tangent
dimension one. Each relative Abel curve $A_i+C$ lies in relative $W^1_5$.
Its differential is injective. Lemma~\ref{lem:formal} identifies the
whole local $W^1_5$ with this known branch.

A hypothetical rational quintic pencil would therefore lie over $\Q_p$
on some $A_i+C$. Upper semicontinuity bounds its number of sections by
two. Both sections then vanish at the added point, contradicting
base-point-freeness. The rank-zero lemma gives finiteness.
The tangent calculation is the exact Petri formula
\[
 \dim T_{[D]}W^1_5=g-\operatorname{rank}
 \bigl(H^0(D)\otimes H^0(K-D)\longrightarrow H^0(K)\bigr)
 \quad(h^0(D)=2).
\]

\section{Positive rank: density theorems and elliptic maps}
Here it is necessary to exclude infinite divisor families, not merely
moving linear systems. First check from the lower-degree classifications
that points of each degree one through four are finite. Quintic density
would then have minimum density degree five. Kadets--Vogt's genus bound
\cite[Theorem 1.3]{KV}
forces a degree-five map to $\PP^1$ or a positive-rank elliptic curve
when $g\ge12$. Their arithmetic classification for $g\ge9$ additionally
allows a Debarre--Fahlaoui normalization
\cite[Proposition 6.3, arithmetic row]{KV}.
All 51 uses of these density theorems have been compared, by full subgroup
element sets, with the lower-degree lists in \cite[Theorems 1.4--1.6]{DO};
the executable \texttt{check\_density\_hypotheses.py} records that check.

The latter alternative lies on $\operatorname{Sym}^2E$ in numerical class
$B=(5+m)H-mF$, $1\le m\le5$, with $H^2=H\cdot F=1$, $F^2=0$,
and $K=-2H+F$. Direct adjunction gives
\[
 p_a(B)=1+\tfrac12(B^2+B\cdot K)=11-\tfrac12m(m-1).
\]
If the normalization has genus at least nine, then $m=1$ or $2$.
The addition map to $E$ consequently has degree $B\cdot F=6$ or $7$.
Thus excluding elliptic maps of degrees five, six and seven handles
these alternatives as well. The arithmetic construction, over $\Q$,
is needed here; the geometric classification alone would not suffice.

At the levels where the optimal-parametrization theorem applies, every
positive-rank elliptic factor has full conductor $N$, multiplicity one,
and odd analytic rank. For $N\le800$ the composition through $X_0(N)$
is optimal \cite[Corollary 4.10]{DO}. If
$t=[(\mathbb Z/N\mathbb Z)^\times:\Delta]$ and $d$
is the optimal modular degree on $X_0(N)$, every elliptic map in that
isogeny class has degree a multiple of $td$. In the applications $t,d>1$,
so degree five is impossible. At levels 37 and 57 the sharper bounds
are 12 and eight. At 73, 109 and 125 the relevant factor computations
exclude positive-rank elliptic quotients. Every application also needs
its separate exclusion of rational quintic functions.

\section{Positive rank: the remaining elliptic families}
Suppose the only potentially positive-rank simple factor of $J(C)$ is
one copy of an elliptic curve. Let $A\subseteq J(C)$ be that elliptic
subvariety. The rational Mordell--Weil group has Zariski closure contained
in finitely many translates of $A$. An infinite intersection of $W_5$
with one such translate is dense in that elliptic curve, so the whole
translate lies in $W_5$.

At level 43, good reduction at two gives $|A(\mathbb F_2)|=5$ and
$J(\mathbb F_2)\simeq\mathbb Z/38\oplus\mathbb Z/1330$. There is
a unique order-five subgroup. Exhaustive enumeration gives
$|W_5(\mathbb F_2)|=2133$ and no full coset of that subgroup in $W_5$.
Prime-to-two torsion specializes injectively, so no rational elliptic
translate can exist. Thus $W_5(\Q)$ is finite.

For levels 53, 61 and the three exceptional level-65 curves, there is
a degree-four map $f:C\to E$ whose pullback image is $A$. The surface
\[
 Z=\{f^*Q+P:Q\in\Pic^1(E),\ P\in C\}\subseteq\Pic^5(C)
\]
is an unavoidable family of reducible effective divisors. At a good odd
prime $p$, let $B=A(\Fp)$. The decisive finite-field tests are:
\begin{enumerate}
\item every full $B$-coset in $W_5(\Fp)$ is a known fiber-plus-point coset;
\item the intersection of the corresponding translated tangent spaces
      has dimension at most two;
\item $A\times C\to\Pic^5(C)$ has rank-two differential at every
      relevant rational reduction preimage.
\end{enumerate}
Since $p\nmid4$, $f_*f^*=[4]$ identifies a smooth elliptic subgroup
over $\mathbb Z_p$. Lift each $b\in B$ to some $a_b\in A(\mathbb Z_p)$
and intersect the relative loci $W_5-a_b$. Neither torsion lifts nor
a subgroup of lifts are required. Both $Z$ and every hypothetical
$A$-translate in $W_5$ are contained in this intersection. The tests and
Lemma~\ref{lem:formal} identify its completion at each relevant point
with the known smooth two-dimensional branch. Consequently each such
elliptic translate is contained in the closed surface $Z$.

An irreducible quintic divisor cannot occur on $Z$ when no quintic
function exists. Indeed, its $h^0$ must be one, so its unique effective
representative is $f^*Q+P$. The leftover point $P$ is unique: distinct
fibers are disjoint, whereas removing two points from a degree-five
divisor leaves common support of degree at least three. Galois invariance
therefore makes $P$ rational, contradicting irreducibility. Only finitely
many isolated rational classes remain outside these families, and each
contributes at most one irreducible divisor.

The tangent computations use a consistent differential basis. At a class
with $h^0(D)=1$, the annihilator of $T_{[D]}W_d$ is $H^0(K-D)$.
Ignoring classes with more sections only weakens the upper bound.
Intersections of tangents at known quartic fibers identify the line
$T_0A$; comparison with Abel tangents checks the final immersion.

The numerical checks for these five curves are summarized below. The
last column counts full $B$-cosets in $W_5(\Fp)$, all of which are known
fiber-plus-point cosets. Every one has tangent-intersection dimension two
and passes the immersion test.
\begin{center}\small
\begin{tabular}{rrrrrr}\toprule
$N$ & Generators of $\Delta$ & $g$ & $p$ & $|B|$ & Full cosets\\\midrule
53 & $\langle-1,4\rangle$ & 8 & 5 & 6 & 8\\
61 & $\langle-1,4\rangle$ & 8 & 5 & 9 & 10\\
65 & $\langle-1,2,7\rangle$ & 9 & 7 & 12 & 4\\
65 & $\langle-1,3\rangle$ & 11 & 3 & 6 & 8\\
65 & $\langle-1,6\rangle$ & 11 & 3 & 6 & 4\\\bottomrule
\end{tabular}
\end{center}

\section{Exceptional case index}
Table~\ref{tab:special} specifies the curves requiring the three
divisor-class methods and the elliptic-family method. Full subgroup
element sets and certificate paths are in \texttt{all\_cases.csv}.
The executable \texttt{replay\_special.py} assigns each of these twelve
cases a separate working copy, re-executes its certificate, and retains
the resulting evidence. It reuses the supplied canonical-model data;
that distinction from fresh model generation is intentional.
\begingroup\small
\begin{table}[!ht]
\centering
\caption{The twelve exceptional negative certificates.}\label{tab:special}
\begin{tabular}{rrrl}
\toprule $N$ & Generators of $\Delta$ & $g$ & Method \\
\midrule
36 & $\langle -1, 17\rangle$ & 7 & Petri lifting \\
39 & $\langle -1, 5\rangle$ & 9 & Petri lifting \\
40 & $\langle -1, 3\rangle$ & 7 & Rational classes \\
43 & $\langle -1, 2\rangle$ & 9 & Empty cosets \\
48 & $\langle -1, 5\rangle$ & 7 & Rational classes \\
48 & $\langle -1, 7, 17\rangle$ & 7 & Rational classes \\
53 & $\langle -1, 4\rangle$ & 8 & Elliptic families \\
61 & $\langle -1, 4\rangle$ & 8 & Elliptic families \\
65 & $\langle -1, 2, 7\rangle$ & 9 & Elliptic families \\
65 & $\langle -1, 3\rangle$ & 11 & Elliptic families \\
65 & $\langle -1, 6\rangle$ & 11 & Elliptic families \\
75 & $\langle -1, 4\rangle$ & 9 & Petri lifting \\
\bottomrule
\end{tabular}
\end{table}

\endgroup

\section{Conclusion and verification scope}
\begin{proof}[Proof of the classification theorem]
The completed $X_0$ classification and the projection lemma exclude all
$N\ge192$. The independent subgroup census lists every remaining pair.
For the 43 strict pairs in Table~\ref{tab:positive}, the indicated
construction gives an exact degree-five function, hence infinitely many
quintic points. Table~\ref{tab:negative} assigns each of 154 strict pairs
one of the direct negative arguments proved above. The remaining negative
pairs are the 807 projections to negative $X_0$ curves and the seven
explicit other forgetting maps in the appendix. Every projection target
has genus greater than one and an already established negative result.
These disjoint assignments account for all $43+154+807+7=1011$ strict
pairs. The endpoint arguments at the start complete the classification.
\end{proof}

The code checks the stated finite-field, divisor-class, rank and model
identities; the deductions from these identities are the mathematical
arguments above. This is not a formal proof-assistant verification.
The command \texttt{python3 verify.py} rebuilds the case assignments from the
saved evidence and checks them against the tables, without Sage or Magma.
The machine-readable \texttt{all\_cases.csv} retains the complete subgroup,
answer, projection target and proof metadata for every enumerated pair.
It is an index to the evidence, not a substitute for the arguments.

\appendix
\section{Complete direct negative assignments}
The table below accounts for all 154 direct negative assignments. Together
with the 807 projections to negative $X_0$ curves and seven other forgetting
maps, these give $154+807+7=968$ negative strict pairs. The 43 positive
pairs complete the 1011 strict pairs through level 191; projection to
$X_0$ excludes all larger levels. The CSV records the exact target for
each forgetting map and the certificate or published input for every row.

For convenience, the seven non-$X_0$ projections are displayed explicitly:
\begin{center}\small
\begin{tabular}{rlrl}
\toprule $N$ & Source generators & $M$ & Target generators \\
\midrule
63 & $\langle -1, 5\rangle$ & 63 & $\langle -1, 4, 5\rangle$ \\
68 & $\langle -1, 13\rangle$ & 34 & $\langle -1, 13\rangle$ \\
68 & $\langle -1, 33\rangle$ & 34 & $\langle -1, 13\rangle$ \\
80 & $\langle -1, 3\rangle$ & 40 & $\langle -1, 3\rangle$ \\
80 & $\langle -1, 9, 31\rangle$ & 40 & $\langle -1, 3\rangle$ \\
80 & $\langle -1, 11\rangle$ & 40 & $\langle -1, 11\rangle$ \\
80 & $\langle -1, 19\rangle$ & 40 & $\langle -1, 19\rangle$ \\
\bottomrule
\end{tabular}

\end{center}

The proof codes refer to the preceding arguments:
\begin{center}\small
\begin{tabular}{ll}\toprule
Code & Argument\\\midrule
S & Spectral gonality bound $119\mu>120000$\\
Z & Rank zero and published rational gonality at least six\\
B & Rank zero and the bielliptic double-cover inequality\\
R & Complete rational divisor-class sieve\\
T & Rank zero and Petri lifting\\
E & Minimum density five and elliptic-degree obstruction\\
E0 & Minimum density five and no positive-rank elliptic factor\\
D & Debarre--Fahlaoui obstruction in genus nine or ten\\
C & No elliptic translate survives reduction\\
F & Surviving elliptic families contain only reducible divisors\\\bottomrule
\end{tabular}
\end{center}
\begingroup\small
\begin{longtable}{rrrrl}
\caption{The $154$ strict pairs with a direct negative argument: the $139$ curves of \S\ref{sec:quintic-intermediate} and the $15$ curves of large index (code S, $119\mu>120000$). Projections are in \cref{tab:projection}.}\label{tab:negative}\\
\toprule $N$ & Generators of $\Delta$ & $g$ & $\mu$ & Proof \\
\midrule\endfirsthead
\toprule $N$ & Generators of $\Delta$ & $g$ & $\mu$ & Proof \\
\midrule\endhead
29 & $\langle -1, 12\rangle$ & 8 & 210 & Z \\
33 & $\langle -1, 10\rangle$ & 11 & 240 & Z \\
34 & $\langle -1, 13\rangle$ & 9 & 216 & Z \\
35 & $\langle -1, 8\rangle$ & 7 & 144 & B \\
35 & $\langle -1, 11\rangle$ & 9 & 192 & Z \\
35 & $\langle -1, 6\rangle$ & 13 & 288 & Z \\
36 & $\langle -1, 17\rangle$ & 7 & 216 & T \\
37 & $\langle -1, 10\rangle$ & 10 & 228 & D \\
37 & $\langle -1, 6\rangle$ & 16 & 342 & E \\
38 & $\langle -1, 7\rangle$ & 10 & 180 & Z \\
39 & $\langle -1, 5\rangle$ & 9 & 168 & T \\
39 & $\langle -1, 16\rangle$ & 9 & 224 & Z \\
39 & $\langle -1, 14\rangle$ & 17 & 336 & Z \\
40 & $\langle -1, 3\rangle$ & 7 & 144 & R \\
40 & $\langle -1, 9\rangle$ & 13 & 288 & Z \\
40 & $\langle -1, 11\rangle$ & 13 & 288 & Z \\
40 & $\langle -1, 19\rangle$ & 9 & 288 & Z \\
41 & $\langle -1, 4\rangle$ & 11 & 168 & Z \\
41 & $\langle -1, 3\rangle$ & 11 & 210 & Z \\
41 & $\langle -1, 9\rangle$ & 21 & 420 & Z \\
42 & $\langle -1, 5\rangle$ & 9 & 192 & Z \\
42 & $\langle -1, 13\rangle$ & 13 & 288 & Z \\
43 & $\langle -1, 2\rangle$ & 9 & 132 & C \\
43 & $\langle -1, 6\rangle$ & 15 & 308 & E \\
44 & $\langle -1, 21\rangle$ & 16 & 360 & Z \\
45 & $\langle -1, 8\rangle$ & 11 & 216 & Z \\
45 & $\langle -1, 14\rangle$ & 9 & 288 & Z \\
45 & $\langle -1, 19\rangle$ & 21 & 432 & Z \\
48 & $\langle -1, 5\rangle$ & 7 & 192 & R \\
48 & $\langle -1, 7, 17\rangle$ & 7 & 192 & R \\
48 & $\langle -1, 7\rangle$ & 19 & 384 & Z \\
48 & $\langle -1, 17\rangle$ & 19 & 384 & Z \\
48 & $\langle -1, 23\rangle$ & 13 & 384 & Z \\
49 & $\langle -1, 18\rangle$ & 19 & 392 & Z \\
50 & $\langle -1, 7\rangle$ & 22 & 450 & Z \\
51 & $\langle -1, 2\rangle$ & 9 & 144 & Z \\
51 & $\langle -1, 4\rangle$ & 17 & 288 & Z \\
51 & $\langle -1, 16\rangle$ & 33 & 576 & Z \\
52 & $\langle -1, 3\rangle$ & 9 & 168 & Z \\
52 & $\langle -1, 5\rangle$ & 13 & 252 & Z \\
52 & $\langle -1, 9\rangle$ & 19 & 336 & Z \\
52 & $\langle -1, 25\rangle$ & 25 & 504 & Z \\
53 & $\langle -1, 4\rangle$ & 8 & 108 & F \\
53 & $\langle -1, 23\rangle$ & 40 & 702 & E \\
54 & $\langle -1, 17\rangle$ & 10 & 324 & Z \\
55 & $\langle -1, 4\rangle$ & 9 & 144 & B \\
55 & $\langle -1, 16\rangle$ & 17 & 288 & Z \\
55 & $\langle -1, 12\rangle$ & 21 & 360 & Z \\
55 & $\langle -1, 21\rangle$ & 41 & 720 & Z \\
56 & $\langle -1, 3\rangle$ & 9 & 192 & Z \\
56 & $\langle -1, 5\rangle$ & 11 & 192 & Z \\
56 & $\langle -1, 17\rangle$ & 11 & 192 & Z \\
56 & $\langle -1, 13, 15\rangle$ & 13 & 288 & Z \\
56 & $\langle -1, 9\rangle$ & 21 & 384 & Z \\
56 & $\langle -1, 13\rangle$ & 31 & 576 & Z \\
56 & $\langle -1, 15\rangle$ & 31 & 576 & Z \\
56 & $\langle -1, 27\rangle$ & 25 & 576 & Z \\
57 & $\langle -1, 2\rangle$ & 9 & 160 & D \\
57 & $\langle -1, 11\rangle$ & 13 & 240 & E \\
57 & $\langle -1, 7\rangle$ & 25 & 480 & E \\
57 & $\langle -1, 20\rangle$ & 43 & 720 & E \\
58 & $\langle -1, 5\rangle$ & 12 & 180 & E \\
58 & $\langle -1, 17\rangle$ & 36 & 630 & E \\
61 & $\langle -1, 4\rangle$ & 8 & 124 & F \\
61 & $\langle -1, 8\rangle$ & 12 & 186 & E \\
61 & $\langle -1, 21\rangle$ & 16 & 310 & E \\
61 & $\langle -1, 3\rangle$ & 26 & 372 & E \\
61 & $\langle -1, 13\rangle$ & 36 & 620 & E \\
61 & $\langle -1, 11\rangle$ & 56 & 930 & E \\
63 & $\langle -1, 4, 5\rangle$ & 9 & 192 & Z \\
63 & $\langle -1, 2\rangle$ & 17 & 288 & Z \\
63 & $\langle -1, 11\rangle$ & 17 & 288 & Z \\
63 & $\langle -1, 10\rangle$ & 17 & 288 & Z \\
63 & $\langle -1, 13\rangle$ & 13 & 288 & Z \\
63 & $\langle -1, 4\rangle$ & 33 & 576 & Z \\
63 & $\langle -1, 17\rangle$ & 33 & 576 & Z \\
63 & $\langle -1, 20\rangle$ & 25 & 576 & Z \\
63 & $\langle -1, 8\rangle$ & 49 & 864 & Z \\
64 & $\langle -1, 15\rangle$ & 13 & 384 & Z \\
64 & $\langle -1, 31\rangle$ & 37 & 768 & Z \\
65 & $\langle -1, 2, 7\rangle$ & 9 & 168 & F \\
65 & $\langle -1, 3\rangle$ & 11 & 168 & F \\
65 & $\langle -1, 6\rangle$ & 11 & 168 & F \\
65 & $\langle -1, 8, 12\rangle$ & 13 & 252 & E \\
65 & $\langle -1, 2\rangle$ & 19 & 336 & E \\
65 & $\langle -1, 7\rangle$ & 19 & 336 & E \\
65 & $\langle -1, 9\rangle$ & 21 & 336 & E \\
65 & $\langle -1, 8, 14\rangle$ & 25 & 504 & E \\
65 & $\langle -1, 12\rangle$ & 31 & 504 & E \\
65 & $\langle -1, 21\rangle$ & 31 & 504 & E \\
65 & $\langle -1, 4\rangle$ & 41 & 672 & E \\
65 & $\langle -1, 8\rangle$ & 55 & 1008 & E \\
65 & $\langle -1, 14\rangle$ & 61 & 1008 & E \\
65 & $\langle -1, 18\rangle$ & 55 & 1008 & E \\
67 & $\langle -1, 3\rangle$ & 15 & 204 & E \\
67 & $\langle -1, 29\rangle$ & 45 & 748 & E \\
68 & $\langle -1, 9\rangle$ & 13 & 216 & Z \\
72 & $\langle -1, 5\rangle$ & 13 & 288 & Z \\
72 & $\langle -1, 7\rangle$ & 13 & 288 & Z \\
72 & $\langle -1, 11\rangle$ & 9 & 288 & Z \\
72 & $\langle -1, 17, 19\rangle$ & 21 & 432 & Z \\
72 & $\langle -1, 23\rangle$ & 25 & 576 & Z \\
72 & $\langle -1, 17\rangle$ & 49 & 864 & Z \\
72 & $\langle -1, 19\rangle$ & 49 & 864 & Z \\
72 & $\langle -1, 35\rangle$ & 41 & 864 & Z \\
73 & $\langle -1, 6\rangle$ & 9 & 148 & D \\
73 & $\langle -1, 7\rangle$ & 13 & 222 & E \\
73 & $\langle -1, 2\rangle$ & 19 & 296 & E \\
73 & $\langle -1, 3\rangle$ & 25 & 444 & E \\
73 & $\langle -1, 10\rangle$ & 43 & 666 & E \\
73 & $\langle -1, 8\rangle$ & 55 & 888 & E \\
73 & $\langle -1, 27\rangle$ & 85 & 1332 & S \\
75 & $\langle -1, 4\rangle$ & 9 & 240 & T \\
75 & $\langle -1, 14\rangle$ & 17 & 480 & Z \\
75 & $\langle -1, 7\rangle$ & 37 & 600 & Z \\
75 & $\langle -1, 26\rangle$ & 73 & 1200 & S \\
80 & $\langle -1, 3, 13\rangle$ & 15 & 288 & Z \\
80 & $\langle -1, 7, 9\rangle$ & 15 & 288 & Z \\
80 & $\langle -1, 9, 11\rangle$ & 13 & 288 & Z \\
80 & $\langle -1, 7\rangle$ & 35 & 576 & Z \\
80 & $\langle -1, 13\rangle$ & 35 & 576 & Z \\
80 & $\langle -1, 17\rangle$ & 35 & 576 & Z \\
80 & $\langle -1, 9\rangle$ & 69 & 1152 & S \\
80 & $\langle -1, 31\rangle$ & 69 & 1152 & S \\
80 & $\langle -1, 39\rangle$ & 57 & 1152 & S \\
81 & $\langle -1, 8\rangle$ & 10 & 324 & Z \\
81 & $\langle -1, 26\rangle$ & 46 & 972 & Z \\
91 & $\langle -1, 3, 4\rangle$ & 13 & 224 & E \\
91 & $\langle -1, 2\rangle$ & 21 & 336 & E \\
91 & $\langle -1, 11\rangle$ & 21 & 336 & E \\
91 & $\langle -1, 5\rangle$ & 23 & 336 & E \\
91 & $\langle -1, 6\rangle$ & 23 & 336 & E \\
91 & $\langle -1, 9, 12\rangle$ & 25 & 448 & E \\
91 & $\langle -1, 3\rangle$ & 41 & 672 & E \\
91 & $\langle -1, 4\rangle$ & 41 & 672 & E \\
91 & $\langle -1, 25\rangle$ & 45 & 672 & E \\
91 & $\langle -1, 36\rangle$ & 45 & 672 & E \\
91 & $\langle -1, 8\rangle$ & 67 & 1008 & E \\
91 & $\langle -1, 9\rangle$ & 81 & 1344 & S \\
91 & $\langle -1, 12\rangle$ & 89 & 1344 & S \\
91 & $\langle -1, 16\rangle$ & 81 & 1344 & S \\
91 & $\langle -1, 22\rangle$ & 89 & 1344 & S \\
91 & $\langle -1, 27\rangle$ & 133 & 2016 & S \\
109 & $\langle -1, 3\rangle$ & 16 & 220 & E0 \\
109 & $\langle -1, 2\rangle$ & 22 & 330 & E0 \\
109 & $\langle -1, 4\rangle$ & 46 & 660 & E0 \\
109 & $\langle -1, 8\rangle$ & 64 & 990 & E0 \\
109 & $\langle -1, 45\rangle$ & 136 & 1980 & S \\
109 & $\langle -1, 33\rangle$ & 208 & 2970 & S \\
121 & $\langle -1, 10\rangle$ & 26 & 660 & E \\
121 & $\langle -1, 3\rangle$ & 106 & 1452 & S \\
125 & $\langle -1, 7\rangle$ & 36 & 750 & E0 \\
125 & $\langle -1, 24\rangle$ & 76 & 1500 & S \\
125 & $\langle -1, 57\rangle$ & 256 & 3750 & S \\
\bottomrule
\end{longtable}

\endgroup

\begin{thebibliography}{99}
\raggedright
\bibitem{DS} M. Derickx and A. V. Sutherland,
\emph{Torsion subgroups of elliptic curves over quintic and sextic number fields},
Proc. Amer. Math. Soc. 145 (2017), 4233--4245.
\url{https://math.mit.edu/~drew/proc13605.pdf}.
\bibitem{Kato} K. Kato,
\emph{$p$-adic Hodge theory and values of zeta functions of modular forms},
Ast\'erisque 295 (2004), 117--290, especially Corollary 14.3.
\url{https://www.numdam.org/item/AST_2004__295__117_0/}.
\bibitem{Orlic} P. Orli\'c, \emph{Tetragonal intermediate modular curves},
arXiv:2407.14512v1, Theorems 1.3--1.5.
\url{https://arxiv.org/html/2407.14512v1}.
\bibitem{KV} B. Kadets and I. Vogt,
\emph{Subspace configurations and low degree points on curves},
arXiv:2208.01067v2.
\url{https://arxiv.org/html/2208.01067v2}.
\bibitem{DO} M. Derickx and P. Orli\'c,
\emph{Intermediate modular curves with infinitely many quartic points},
arXiv:2504.15937v1.
\url{https://arxiv.org/html/2504.15937v1}.
\bibitem{JKS} D. Jeon, C. H. Kim and A. Schweizer,
\emph{Bielliptic intermediate modular curves}, arXiv:1709.00510.
\url{https://arxiv.org/html/1709.00510}.
\bibitem{DEH} M. Derickx, A. Etropolski, M. van Hoeij, J. S. Morrow
and D. Zureick-Brown, \emph{Sporadic cubic torsion}, arXiv:2007.13929v2.
\url{https://arxiv.org/html/2007.13929v2}.
\bibitem{DHJO} M. Derickx, W. Hwang, D. Jeon and P. Orli\'c,
\emph{Modular curves $X_0(N)$ of density degree 5}, arXiv:2503.08975v2.
\url{https://arxiv.org/html/2503.08975v2}.
\bibitem{CES} B. Conrad, B. Edixhoven and W. Stein,
\emph{$J_1(p)$ has connected fibers}, Documenta Math. 8 (2003), 325--402.
\url{https://virtualmath1.stanford.edu/~conrad/papers/j1p.pdf}.
\bibitem{Faltings} G. Faltings,
\emph{Diophantine approximation on abelian varieties},
Ann. of Math. 133 (1991), 549--576.
\url{https://annals.math.princeton.edu/1991/133-3/p04}.
\end{thebibliography}
\end{document}
